%% 1-16 1より 小さい かず：小数と 分数（10ページ、p.160〜169）第3部の さいごの 章
%% 第1巻のきまり：15 pt・行送り28 pt、分かち書き（文節の間に半角の空白）、漢字はすべてルビ、式は全角の ＋ ＝ で書く

%% ---------------- この章だけの絵の部品（\bkqo…） ----------------
\makeatletter
% かんがえかたの 2行めを、 1行めの 書きだしに そろえる 字下げ
\newcommand\bkqokindent{\setbox0=\hbox{\tikz\node[font=\sffamily\bfseries\small, inner xsep=5pt, inner ysep=2.5pt]{かんがえかた};}\hspace*{\dimexpr\wd0+0.6\zw\relax}}
% 分数（数字は本文と同じ字）  \bkqofrac{分子}{分母}
\newcommand\bkqofrac[2]{\mbox{\(\dfrac{\mbox{#1}}{\mbox{#2}}\)}}
% dL ます（目もり10）  \bkqomasu{(x,y)}{はば}{高さ}{水の目もり(0〜10)}   (x,y) は底のまん中。ふちが ちょうど 1 dL
\newcommand\bkqomasu[4]{\begin{scope}[shift={#1}]
  \ifnum#4>0 \fill[bkBlue!30] (-#2/2+0.55,0.55) rectangle (#2/2-0.55,{#3*#4/10});\fi
  \foreach \k in {1,...,9} {\pgfmathsetmacro\bkqo@t{(\k==5)?4.6:2.8}\draw[bkInk!75, line width=0.5pt] (-#2/2,{#3*\k/10}) -- ++(\bkqo@t,0);}
  \draw[bkInk!80, line width=1.1pt, line join=round] (-#2/2-0.8,#3) -- (-#2/2,#3) -- (-#2/2,0) -- (#2/2,0) -- (#2/2,#3) -- (#2/2+0.8,#3);
\end{scope}}
% ものさし（じっさいの 大きさ：1 cm＝10 mm）  \bkqoruler{(x,y)}{cmの数}   (x,y) は 0の目もりの上のはし。目もりは下へ
\newcommand\bkqoruler[2]{\begin{scope}[shift={#1}]
  \pgfmathtruncatemacro\bkqo@n{#2*10}
  \fill[bkYellowBg, draw=bkYellow!80!black, line width=0.7pt, rounded corners=0.8mm] (-3,0) rectangle ({#2*10+3},-10);
  \foreach \m in {0,...,\bkqo@n} {\pgfmathsetmacro\bkqo@t{(mod(\m,10)==0)?4.4:((mod(\m,5)==0)?3.2:2)}
    \draw[bkInk, line width=0.35pt] (\m,0) -- (\m,-\bkqo@t);}
  \foreach \c in {0,...,#2} {\node[font=\sffamily\bfseries\fontsize{8}{9}\selectfont, text=bkInk, inner sep=0pt] at ({\c*10},-7.2) {\c};}
\end{scope}}
% 小数の 数直線  \bkqonum{(x,y)}{1の長さ}{さいごの整数}
\newcommand\bkqonum[3]{\begin{scope}[shift={#1}]
  \pgfmathtruncatemacro\bkqo@n{#3*10}
  \draw[bkInk, line width=0.9pt] (-2,0) -- ({#3*#2+3},0);
  \foreach \k in {0,...,\bkqo@n} {\pgfmathsetmacro\bkqo@t{(mod(\k,10)==0)?3:((mod(\k,5)==0)?2.2:1.4)}
    \draw[bkInk, line width=0.6pt] ({\k*#2/10},-\bkqo@t) -- ({\k*#2/10},\bkqo@t);}
  \foreach \c in {0,...,#3} {\node[anchor=north, font=\sffamily\bfseries\small, text=bkInk, inner sep=0pt] at ({\c*#2},-4) {\c};}
\end{scope}}
% 数直線の上の 矢じるし  \bkqomark{(x,y)}{文字}{色}
\newcommand\bkqomark[3]{\begin{scope}[shift={#1}]
  \draw[#3, line width=1.2pt, -{Stealth[length=4pt]}] (0,8.2) -- (0,3.4);
  \node[anchor=east, font=\sffamily\bfseries\small, text=#3, inner sep=1pt] at (-0.4,7) {#2};
\end{scope}}
% 同じ長さに 分けた テープ  \bkqotape{(x,y)}{長さ}{分ける数}{ぬる数}{色}
\newcommand\bkqotape[5]{\begin{scope}[shift={#1}]
  \ifnum#4>0 \fill[#5] (0,0) rectangle ({#2*#4/#3},7);\fi
  \draw[bkInk!80, line width=0.8pt] (0,0) rectangle (#2,7);
  \ifnum#3>1 \foreach \k in {1,...,\numexpr#3-1\relax} {\draw[bkInk!80, line width=0.8pt] ({#2*\k/#3},0) -- ({#2*\k/#3},7);}\fi
\end{scope}}
% ぐの のった ピザ（はてな？用）  \bkqopizza{(x,y)}{半径}{分ける数(1・2・4)}
\newcommand\bkqopizza[3]{\begin{scope}[shift={#1}]
  \fill[bkStem!45] (0,0) circle (#2);
  \fill[bkYellow!40] (0,0) circle ({#2*0.86});
  \foreach \a/\d in {25/0.55,110/0.52,205/0.58,300/0.5,160/0.22,345/0.25} {\fill[bkAppleRed!80] (\a:{#2*\d}) circle ({#2*0.12});}
  \foreach \a/\d in {65/0.66,175/0.66,250/0.7,330/0.68,10/0.8} {\fill[bkLeaf] (\a:{#2*\d}) circle ({#2*0.065});}
  \ifnum#3>1 \foreach \k in {1,...,#3} {\draw[white, line width=1.6pt] (0,0) -- ({90+360*\k/#3}:#2);}\fi
\end{scope}}
\makeatother

\bkchapter{1}{16}{1より \R{小}{ちい}さい かず：\\\mbox{\R{小数}{しょう|すう}}と \mbox{\R{分数}{ぶん|すう}}}{\bklevel{1}}{1より ちいさい かず：しょうすうと ぶんすう}
\bkmichishirube{1-6,1-8,1-9}{2-1,2-4,2-6}

\begin{goal}
\begin{itemize}[label={\color{bkTheme}●}, leftmargin=1.4\zw, itemsep=0pt, topsep=0pt]
\item はんぱな かさや \R{長}{なが}さを、 \mbox{\R{小数}{しょう|すう}}で あらわせる。
\item \mbox{\R{数直線}{すう|ちょく|せん}}で、 \mbox{\R{小数}{しょう|すう}}の \R{目}{め}もりが よめる。
\item \mbox{\R{半分}{はん|ぶん}}や \mbox{\R{半分}{はん|ぶん}}の \mbox{\R{半分}{はん|ぶん}}を、 \mbox{\R{分数}{ぶん|すう}}で あらわせる。
\end{itemize}
\end{goal}

\begin{hatena}
\begin{bkfigure}
\begin{tikzpicture}[x=1mm, y=1mm]
  \fill[white, rounded corners=3mm] (0,0) rectangle (122,48);
  \bkqopizza{(18,25)}{15}{1}
  \bkqopizza{(61,25)}{15}{2}
  \bkqopizza{(104,25)}{15}{4}
  \draw[bkTheme, line width=1.8pt, line join=round] (104,25) -- (104,40) arc (90:0:15) -- cycle;
  \foreach \x in {35,78} {\draw[bk arrow, bkTheme, line width=1.6pt] (\x,25) -- ({\x+9},25);}
  \node[font=\sffamily\bfseries\small, text=bkInk] at (18,4.5) {1まい};
  \node[font=\sffamily\bfseries\small, text=bkInk] at (61,4.5) {\R{半分}{はん|ぶん}に};
  \node[font=\sffamily\bfseries\small, text=bkInk] at (104,4.5) {また \mbox{\R{半分}{はん|ぶん}}に};
  \node[circle, fill=bkYellow, text=white, font=\sffamily\bfseries\Large, minimum size=8mm, inner sep=0pt] at (117,42) {？};
\end{tikzpicture}
\end{bkfigure}
\noindent ピザ 1まいを、 \R{同}{おな}じ \R{大}{おお}きさに \mbox{\R{半分}{はん|ぶん}}に \R{切}{き}りました。 それを また \mbox{\R{半分}{はん|ぶん}}に \R{切}{き}りました。\\
この 1きれは、 ピザ 1まいの どれだけ かな？ 「1まい」より \R{小}{ちい}さいね。
\end{hatena}

\clearpage
%% ---------------- p.161 やってみよう ----------------
\begin{yattemiyou}[はんぱを はかって みよう]
\begin{enumerate}[label={\bkcircled[bkGreen]{\arabic*}}, leftmargin=2.2\zw, itemsep=2pt, topsep=0pt]
\item \R{水}{すい}とうの \R{水}{みず}を、 1\,dL ますで はかろう。 2はい はかったら、 すこし のこりました。
\end{enumerate}
\begin{bkfigure}
\begin{tikzpicture}[x=1mm, y=1mm]
  \foreach \x in {20,48} {\bkqomasu{(\x,0)}{20}{26}{10}\node[font=\sffamily\bfseries\small, text=bkInk] at (\x,-4.5) {1\,dL};}
  \bkqomasu{(84,0)}{20}{26}{3}
  \node[font=\sffamily\bfseries\small, text=bkGreen!70!black] at (84,-4.5) {のこり};
  \draw[bkGreen, line width=1pt, -{Stealth[length=4pt]}] (103,14) -- (95.5,8.5);
  \node[anchor=south west, font=\sffamily\bfseries\footnotesize, text=bkGreen!70!black, align=left] at (99,14.5) {\R{目}{め}もり\\いくつ？};
\end{tikzpicture}
\end{bkfigure}
\begin{enumerate}[label={\bkcircled[bkGreen]{\arabic*}}, leftmargin=2.2\zw, itemsep=2pt, topsep=4pt, start=2]
\item のこりの \R{水}{みず}は、 \R{小}{ちい}さい \R{目}{め}もり いくつ\R{分}{ぶん}かな？
\end{enumerate}
\par\vspace{2pt}
\noindent\hspace*{2.2\zw}1\,dL が\enspace\bkkakikomi\enspace はいと、 \R{目}{め}もり\enspace\bkkakikomi\enspace こ\R{分}{ぶん}
\begin{enumerate}[label={\bkcircled[bkGreen]{\arabic*}}, leftmargin=2.2\zw, itemsep=2pt, topsep=6pt, start=3]
\item \R{下}{した}の けしゴムの \R{長}{なが}さを、 \mbox{\R{自分}{じ|ぶん}}の ものさしで はかろう。
\end{enumerate}
\begin{bkfigure}
\begin{tikzpicture}[x=1mm, y=1mm]
  % けしゴム（じっさいの 大きさ：4 cm 3 mm。たかさは 本物に 近い 約 2.3：1）
  \bkeraser{(0,0)}{43}{18.5}% 第11回：全巻共通の消しゴム
  \node[anchor=west, font=\sffamily\bfseries\footnotesize, text=bkInk] at (47,9.25) {けしゴム};
\end{tikzpicture}
\end{bkfigure}
\noindent\hspace*{2.2\zw}\bkkakikomi\enspace cm と\enspace\bkkakikomi\enspace mm
\end{yattemiyou}

\bkfukidashi{mina}{どちらも 「1」に たりない はんぱが あるね。 はんぱを 1つの \R{数}{かず}で いえたら いいのに。}

\clearpage
%% ---------------- p.162〜163 数学の道具 ----------------
\bksection{doguu}{\R{数学}{すう|がく}の \mbox{\R{道具}{どう|ぐ}}}

\bksubsection{1}{1を 10に \R{分}{わ}けた 1つ\R{分}{ぶん}：0.1}

\label{bk:1-16:shousuu}1\,dL を 10に \R{分}{わ}けた 1つ\R{分}{ぶん}を 0.1\,dL と かき、 「れい てん いち デシリットル」と よみます。

\begin{bkfigure}
\begin{tikzpicture}[x=1mm, y=1mm]
  \bkqomasu{(18,0)}{20}{30}{1}
  \draw[bkTheme, line width=1pt] (29,0.3) -- (31,0.3) -- (31,2.7) -- (29,2.7);
  \node[anchor=west, font=\sffamily\bfseries\small, text=bkTheme] at (31.5,1.5) {0.1\,dL};
  \bkbrace{(5,0)}{(5,30)}{}{1}
  \node[anchor=east, font=\sffamily\bfseries\footnotesize, text=bkInk, align=right] at (2,15) {10こで\\1\,dL};
  % やってみよう の 水
  \foreach \x in {62,85} {\bkqomasu{(\x,0)}{20}{30}{10}}
  \bkqomasu{(108,0)}{20}{30}{3}
  \node[font=\sffamily\bfseries\small, text=bkInk] at (73.5,-4.5) {2\,dL};
  \node[font=\sffamily\bfseries\small, text=bkInk] at (108,-4.5) {0.3\,dL};
\end{tikzpicture}
\end{bkfigure}
\noindent なぜ 10に \R{分}{わ}けるのかな？ 1-2で、 1が 10こ あつまると 10に なったね。 こんどは \R{反対}{はん|たい}に、 1を 10こに \R{分}{わ}けた 1つ\R{分}{ぶん}を つくるのです。 だから 0.1が 10こ あつまると、 ちょうど 1に もどります。

\noindent やってみよう の \R{水}{みず}は、 2\,dL と 0.3\,dL で 2.3\,dL（に てん さん デシリットル）。 0.3\,dL は、 0.1\,dL の \R{目}{め}もりが 3こ\R{分}{ぶん}です。 2.3 や 0.1 の ような \R{数}{かず}を 「\mbox{\R{小数}{しょう|すう}}」、 「.」を 「\mbox{\R{小数点}{しょう|すう|てん}}」と いいます。\bkidx{しょうすう}{\R{小数}{しょう|すう}}\bkidx{しょうすうてん}{\R{小数点}{しょう|すう|てん}}

\begin{bkfigure}
\begin{tikzpicture}[x=1mm, y=1mm]
  \tikzset{cl/.style={draw=bkInk!60, line width=0.8pt, minimum width=24mm, minimum height=9mm, font=\sffamily\bfseries\small}}
  \node[cl, fill=bkTheme!10] at (0,9) {\R{一}{いち}の くらい};
  \node[cl, fill=bkBlue!12, minimum width=10mm] at (17,9) {};
  \node[cl, fill=bkTheme!10] at (34,9) {0.1の くらい};
  \node[cl, font=\fontsize{22}{24}\selectfont] at (0,0) {2};
  \node[cl, minimum width=10mm, font=\fontsize{22}{24}\selectfont, text=bkRed] at (17,0) {.};
  \node[cl, font=\fontsize{22}{24}\selectfont] at (34,0) {3};
  \node[anchor=west, font=\sffamily\bfseries\footnotesize, text=bkInk, align=left] at (52,4.5) {1が 2こと、 0.1が 3こ。\\0.1が ぜんぶで 23こ とも いえる。};
\end{tikzpicture}
\end{bkfigure}
\noindent \mbox{\R{小数点}{しょう|すう|てん}}の \R{右}{みぎ}に ある \R{数}{かず}は、 「0.1が いくつ」を あらわします。 1-2の くらいの へやが、 \R{右}{みぎ}に 1つ ふえたのです。

\bksubsection{2}{ものさしの 0.1\,cm}

\label{bk:1-16:monosashi}1\,cm を 10に \R{分}{わ}けた 1つ\R{分}{ぶん}が 1\,mm（1-8）。 だから \mbox{1\,mm ＝ 0.1\,cm} です。

\begin{bkfigure}
\begin{tikzpicture}[x=1mm, y=1mm]
  \bkeraser{(0,0.6)}{43}{18.5}% 第11回：全巻共通の消しゴム（左はしを 0 の 目もりに あわせる）
  \bkqoruler{(0,0)}{6}
  \draw[bkTheme, line width=0.8pt] (43,0.4) -- (43,-4.6);
  \node[anchor=west, font=\sffamily\bfseries\small, text=bkTheme, align=left] at (70,-1) {4\,cm 3\,mm\\＝ 4.3\,cm};
\end{tikzpicture}
\end{bkfigure}
\vspace{2.5mm}
\noindent けしゴムは 4\,cm と、 1\,mm が 3こ。 1\,mm は 0.1\,cm なので、 3\,mm は 0.1\,cm が 3こで 0.3\,cm。 あわせて 4.3\,cm です。

\begin{taisetsu}
1を 10に \R{分}{わ}けた 1つ\R{分}{ぶん}が 0.1。 0.1が 10こで 1。
\end{taisetsu}

\bksubsection{3}{\R{数直線}{すう|ちょく|せん}の \mbox{\R{小数}{しょう|すう}}}

\label{bk:1-16:suuchokusen}\mbox{\R{数直線}{すう|ちょく|せん}}（1-6）の 0と 1の あいだを 10に \R{分}{わ}けると、 \mbox{1\R{目}{め}}もりは 0.1です。

\begin{bkfigure}
\begin{tikzpicture}[x=1mm, y=1mm]
  \bkqonum{(0,0)}{55}{2}
  \bkqomark{(16.5,0)}{0.3}{bkBlue}
  \bkqomark{(71.5,0)}{1.3}{bkTheme}
  \bkbrace{(0,-9.5)}{(55,-9.5)}{0.1が 10こで 1}{-1}
\end{tikzpicture}
\end{bkfigure}
\noindent 0.3は、 0から \R{目}{め}もり 3つ\R{分}{ぶん}。 1.3は、 1から さらに \R{目}{め}もり 3つ\R{分}{ぶん} \R{右}{みぎ}。 0から かぞえると、 0.1が 13こ\R{分}{ぶん}です。 \R{数直線}{すう|ちょく|せん}では、 \R{右}{みぎ}に ある \R{数}{かず}ほど \R{大}{おお}きく なります（1-6）。 だから 1.3は 0.3より \R{大}{おお}きい \R{数}{かず}です。

\bksubsection{4}{\R{半分}{はん|ぶん}と、 \mbox{\R{半分}{はん|ぶん}}の \mbox{\R{半分}{はん|ぶん}}：\mbox{\R{分数}{ぶん|すう}}}

\label{bk:1-16:bunsuu}\R{同}{おな}じ \R{大}{おお}きさに 2つに \R{分}{わ}けた 1つ\R{分}{ぶん}を \bkqofrac{1}{2}（にぶんの いち）、 4つに \R{分}{わ}けた 1つ\R{分}{ぶん}を \bkqofrac{1}{4}（よんぶんの いち）と かきます。\bkidx{ぶんすう}{\R{分数}{ぶん|すう}}

\begin{bkfigure}
\begin{tikzpicture}[x=1mm, y=1mm]
  \fill[bkYellowBg] (14,0) circle (11); \draw[bkYellow!60!black, line width=1.3pt] (14,0) circle (11);
  \bkpizza{(57,0)}{11}{2}{1}{bkOrange}
  \bkpizza{(102,0)}{11}{4}{1}{bkOrange}
  \foreach \x in {29.5,72.5} {\draw[bk arrow, bkTheme, line width=1.4pt] (\x,0) -- ({\x+10},0);}
  \node[font=\sffamily\bfseries\footnotesize, text=bkTheme] at (34.5,4.5) {\R{半分}{はん|ぶん}};
  \node[font=\sffamily\bfseries\footnotesize, text=bkTheme] at (77.5,4.5) {\R{半分}{はん|ぶん}};
  \node[anchor=west, font=\sffamily\bfseries\small] at (26,-9) {1まい};
  \node[anchor=west, font=\sffamily\bfseries\small] at (69,-8) {\bkqofrac{1}{2}};
  \node[anchor=west, font=\sffamily\bfseries\small] at (114,-8) {\bkqofrac{1}{4}};
\end{tikzpicture}
\end{bkfigure}
\noindent\begin{minipage}[c]{0.55\linewidth}
\R{分数}{ぶん|すう}の \R{下}{した}の \R{数}{かず}は、 \R{同}{おな}じ \R{大}{おお}きさに いくつに \R{分}{わ}けたか。 \R{上}{うえ}の \R{数}{かず}は、 その うちの いくつ\R{分}{ぶん}か を あらわします。
\end{minipage}\hfill
\begin{minipage}[c]{0.43\linewidth}\centering
\begin{tikzpicture}[x=1mm, y=1mm]
  \node[font=\fontsize{28}{30}\selectfont] (f) at (0,0) {\(\dfrac{\mbox{1}}{\mbox{4}}\)};
  \draw[bkOrange, line width=0.8pt] (4.5,5.2) -- (10,7.5);
  \node[anchor=west, font=\sffamily\bfseries\footnotesize, text=bkOrange!85!black] at (10.5,7.5) {その うちの 1つ\R{分}{ぶん}};
  \draw[bkBlue, line width=0.8pt] (4.5,-5.2) -- (10,-7.5);
  \node[anchor=west, font=\sffamily\bfseries\footnotesize, text=bkBlue!85!black] at (10.5,-7.5) {4つに \R{分}{わ}けた};
\end{tikzpicture}
\end{minipage}

\noindent \bkhatenaref の ピザは、 \R{半分}{はん|ぶん}に \R{切}{き}って、 それぞれを また \R{半分}{はん|ぶん}に \R{切}{き}ったので、 \R{同}{おな}じ \R{大}{おお}きさの 4きれに なりました。 その 1きれだから、 こたえは ピザ 1まいの \bkqofrac{1}{4} です。

\begin{bkfigure}
\begin{tikzpicture}[x=1mm, y=1mm]
  \fill[bkOrange!55] (0,0) rectangle (24,7); \fill[bkOrange!20] (24,0) rectangle (48,7);
  \draw[bkInk!80, line width=0.8pt] (0,0) rectangle (48,7) (24,0) -- (24,7);
  \node[anchor=west, font=\sffamily\bfseries\footnotesize] at (51,3.5) {\bkqofrac{1}{2} が 2つで 1};
  \begin{scope}[shift={(0,-14)}]
  \foreach \k/\c in {0/55,1/20,2/55,3/20} {\fill[bkBlue!\c] ({\k*12},0) rectangle ++(12,7);}
  \draw[bkInk!80, line width=0.8pt] (0,0) rectangle (48,7) (12,0) -- (12,7) (24,0) -- (24,7) (36,0) -- (36,7);
  \node[anchor=west, font=\sffamily\bfseries\footnotesize] at (51,3.5) {\bkqofrac{1}{4} が 4つで 1};
  \end{scope}
  \draw[bkTheme, line width=1pt, dash pattern=on 2pt off 1.5pt] (0,-16) rectangle (24,9);
  \node[anchor=west, font=\sffamily\bfseries\footnotesize, text=bkTheme] at (51,-21) {\bkqofrac{1}{4} が 2つで \bkqofrac{1}{2} と \R{同}{おな}じ \R{長}{なが}さ};
\end{tikzpicture}
\end{bkfigure}

\bksubsection{5}{かずの \bkqofrac{1}{2}・\bkqofrac{1}{4}}

\label{bk:1-16:kazu}ピザや テープだけで なく、 あめの ような 「こ」で かぞえる ものも、 \R{同}{おな}じ かずずつ \R{分}{わ}けられます。 あめ 8この \bkqofrac{1}{2} は、 8こを \R{同}{おな}じ かずずつ 2つに \R{分}{わ}けた 1つ\R{分}{ぶん}です。
\begin{bkfigure}
\begin{tikzpicture}[x=1mm, y=1mm]
  \foreach \i in {0,...,3} {\bkcandy{({6+\i*8},4)}{1.5}{bkAppleRed}}
  \foreach \i in {0,...,3} {\bkcandy{({46+\i*8},4)}{1.5}{bkRoute2}}
  \draw[bkTheme, line width=1pt, rounded corners=2mm] (0,-2) rectangle (36,10);
  \draw[bkGray, line width=0.8pt, rounded corners=2mm] (40,-2) rectangle (76,10);
  \node[font=\sffamily\bfseries\footnotesize, text=bkTheme] at (18,-6.5) {\bkqofrac{1}{2} は 4こ};
  \node[anchor=west, font=\sffamily\bfseries\small, text=bkInk] at (82,4) {8 ÷ 2 ＝ 4};
\end{tikzpicture}
\end{bkfigure}
\noindent 2つに \R{同}{おな}じ かずずつ \R{分}{わ}けるのは、 1-14の わり\R{算}{ざん} \mbox{8 ÷ 2 ＝ 4} です。 \R{同}{おな}じ ように、 8この \bkqofrac{1}{4} は \mbox{8 ÷ 4 ＝ 2} で 2こ です。

\bksubsection{6}{0.5と \bkqofrac{1}{2} は \R{同}{おな}じ \R{大}{おお}きさ}

\label{bk:1-16:onaji}0から 1までを 10に \R{分}{わ}けた \R{数直線}{すう|ちょく|せん}と、 2つに \R{分}{わ}けた テープを ならべて みよう。
\begin{bkfigure}
\begin{tikzpicture}[x=1mm, y=1mm]
  \bkqotape{(0,10)}{100}{2}{1}{bkOrange!45}
  \node[anchor=south, font=\sffamily\bfseries\small] at (25,17.5) {\bkqofrac{1}{2}};
  \bkqonum{(0,0)}{100}{1}
  \draw[bkTheme, line width=1pt, dash pattern=on 2pt off 1.5pt] (50,10) -- (50,3.2);
  \node[anchor=north, font=\sffamily\bfseries\small, text=bkTheme] at (50,-4) {0.5};
\end{tikzpicture}
\end{bkfigure}
\noindent 1の \R{半分}{はん|ぶん}の ところは、 0.1が 5こ\R{分}{ぶん}で 0.5です。 0.5と \bkqofrac{1}{2} は、 \R{同}{おな}じ \R{大}{おお}きさを ちがう かきかたで あらわして いるのです。 （0.1が 10こで 1。 その \R{半分}{はん|ぶん}は 5こ だからです。）

\bkkigocards{\bkkigocard{2{\color{bkRed}.}3}{に てん さん}{2と 0.3を あわせた \R{数}{かず}。 「.」は \mbox{\R{小数点}{しょう|すう|てん}}}{decimal point}{1-16}}{\bkkigocard{\bkqofrac{1}{2}}{にぶんの いち}{\R{同}{おな}じ \mbox{\R{大}{おお}きさ}に 2つに \R{分}{わ}けた 1つ\R{分}{ぶん}}{one half}{1-16}}

%% ---------------- p.164〜165 例題 ----------------

\begin{reidai}[かさを \mbox{\R{小数}{しょう|すう}}で]
\noindent\begin{minipage}[c]{0.72\linewidth}
\R{水}{みず}の かさは \R{何}{なん}\,dL？

\kangaekata 1\,dL と、 0.1\,dL が 4こ。

\kotae 1.4\,dL
\end{minipage}\hfill
\begin{minipage}[c]{0.26\linewidth}\centering
\begin{tikzpicture}[x=1mm, y=1mm]
  \bkqomasu{(0,0)}{12}{18}{10}\bkqomasu{(15,0)}{12}{18}{4}
\end{tikzpicture}
\end{minipage}
\end{reidai}

\begin{jibun}
\noindent\begin{minipage}[c]{0.6\linewidth}
\R{水}{みず}の かさは \R{何}{なん}\,dL？\par\vspace{4pt}
\hspace*{\fill}{\sffamily\bfseries\small\color{bkOrange}\bkkotaelabel}\enspace\bkkakikomi[24mm]
\bkjans{2.7\,dL}
\end{minipage}\hfill
\begin{minipage}[c]{0.38\linewidth}\centering
\begin{tikzpicture}[x=1mm, y=1mm]
  \bkqomasu{(0,0)}{12}{18}{10}\bkqomasu{(15,0)}{12}{18}{10}\bkqomasu{(30,0)}{12}{18}{7}
\end{tikzpicture}
\end{minipage}
\end{jibun}

\begin{reidai}[\R{長}{なが}さを \mbox{\R{小数}{しょう|すう}}で]
テープの \R{長}{なが}さは \R{何}{なん}\,cm？
\begin{bkfigure}
\begin{tikzpicture}[x=1mm, y=1mm]
  \bktape{(0,1)}{46}{bkRoute6}{}
  \bkqoruler{(0,0)}{6}
\end{tikzpicture}
\end{bkfigure}
\kangaekata 4\,cm と 6\,mm。 6\,mm は 0.6\,cm。 \kotaenext 4.6\,cm
\end{reidai}

\begin{jibun}
えんぴつの \R{長}{なが}さは \R{何}{なん}\,cm？\hfill{\sffamily\bfseries\small\color{bkOrange}\bkkotaelabel}\enspace\bkkakikomi[24mm]
\begin{bkfigure}
\begin{tikzpicture}[x=1mm, y=1mm]
  \bkpencil{(2.2,4)}{49.8}{bkRoute4}
  \bkqoruler{(0,0)}{6}
\end{tikzpicture}
\end{bkfigure}
\bkjans{5.2\,cm}
\end{jibun}

\begin{reidai}[\mbox{\R{数直線}{すう|ちょく|せん}}を よもう]
ア、 イの \R{目}{め}もりは いくつ？
\begin{bkfigure}
\begin{tikzpicture}[x=1mm, y=1mm]
  \bkqonum{(0,0)}{55}{2}
  \bkqomark{(38.5,0)}{ア}{bkBlue}
  \bkqomark{(71.5,0)}{イ}{bkBlue}
\end{tikzpicture}
\end{bkfigure}
\kangaekata \mbox{1\R{目}{め}}もりは 0.1。 イは 0.1が 13こ。

\kotae ア 0.7、 イ 1.3
\end{reidai}

\begin{jibun}
ウ、 エの \R{目}{め}もりは いくつ？\hfill ウ\enspace\bkkakikomi[14mm]\enspace エ\enspace\bkkakikomi[14mm]
\begin{bkfigure}
\begin{tikzpicture}[x=1mm, y=1mm]
  \bkqonum{(0,0)}{55}{2}
  \bkqomark{(22,0)}{ウ}{bkOrange}
  \bkqomark{(99,0)}{エ}{bkOrange}
\end{tikzpicture}
\end{bkfigure}
\bkjans{\mbox{ウ 0.4}、 \mbox{エ 1.8}}
\end{jibun}

\begin{reidai}[どちらが \R{大}{おお}きい？]
0.8 と 1.2 は、 どちらが \R{大}{おお}きい？
\begin{bkfigure}
\begin{tikzpicture}[x=1mm, y=1mm]
  \bkqonum{(0,0)}{55}{2}
  \bkqomark{(44,0)}{0.8}{bkBlue}
  \bkqomark{(66,0)}{1.2}{bkTheme}
\end{tikzpicture}
\end{bkfigure}
\kangaekata \R{数直線}{すう|ちょく|せん}で \R{右}{みぎ}に あるのは 1.2。 0.8は 1より \R{小}{ちい}さく、 1.2は 1より \R{大}{おお}きい。

\kotae 1.2
\end{reidai}

\begin{jibun}
\R{大}{おお}きい ほうに ○を つけよう。\par\vspace{2pt}
\noindent\hspace*{1\zw}（1）1.5\enspace と\enspace 0.9\hspace{3\zw}（2）2.1\enspace と\enspace 1.9
\bkjans{（1）1.5\quad（2）2.1}
\end{jibun}

\bkfukidashi{sora}{1より \R{小}{ちい}さいか、 1より \R{大}{おお}きいかを まず \R{見}{み}ると、 くらべやすいね。 \R{小数点}{しょう|すう|てん}の \R{左}{ひだり}の \R{数}{かず}から くらべよう。}

\begin{reidai}[\mbox{\R{分数}{ぶん|すう}}で あらわそう]
\R{色}{いろ}の ところは、 もとの \R{大}{おお}きさの どれだけ？\par
\noindent\begin{minipage}[c]{0.62\linewidth}
\kangaekata \R{下}{した}は、 \R{分}{わ}けた \R{数}{かず}。

\kotae A\enspace\bkqofrac{1}{4}\quad B\enspace\bkqofrac{1}{2}
\end{minipage}\hfill
\begin{minipage}[c]{0.36\linewidth}\centering
\begin{tikzpicture}[x=1mm, y=1mm]
  \node[circle, fill=bkBlue, text=white, font=\sffamily\bfseries\footnotesize, minimum size=4.5mm, inner sep=0pt] at (-1.5,7) {A};
  \bkpizza{(7,0)}{7.5}{4}{1}{bkBlue!60}
  \node[circle, fill=bkBlue, text=white, font=\sffamily\bfseries\footnotesize, minimum size=4.5mm, inner sep=0pt] at (18,7) {B};
  \bkqotape{(18,-3.5)}{22}{2}{1}{bkBlue!60}
\end{tikzpicture}
\end{minipage}
\end{reidai}

\begin{jibun}
\noindent\begin{minipage}[c]{0.8\linewidth}
おりがみを \mbox{\R{半分}{はん|ぶん}}の \mbox{\R{半分}{はん|ぶん}}に おりました。 \R{色}{いろ}の ところは どれだけ？\hfill{\sffamily\bfseries\small\color{bkOrange}\bkkotaelabel}\enspace\bkkakikomi[16mm]
\bkjans{\bkqofrac{1}{4}}
\end{minipage}\hfill
\begin{minipage}[c]{0.18\linewidth}\centering
\begin{tikzpicture}[x=1mm, y=1mm]
  \fill[bkOrange!15] (0,0) rectangle (18,18);
  \fill[bkOrange!70] (0,9) rectangle (9,18);
  \draw[bkOrange!80!black, line width=0.8pt] (0,0) rectangle (18,18);
  \draw[bkOrange!80!black, line width=0.8pt, dash pattern=on 2pt off 1.5pt] (9,0) -- (9,18) (0,9) -- (18,9);
\end{tikzpicture}
\end{minipage}
\end{jibun}

\begin{reidai}[L を \mbox{\R{小数}{しょう|すう}}で]
\noindent\begin{minipage}[c]{0.66\linewidth}
\R{水}{みず}は \R{何}{なん}\,L？ 1\,Lますの \R{目}{め}もりは、 1\,dL ずつです。

\kangaekata 1\,L は 10\,dL（1-9）。 だから \mbox{1\,dL ＝ 0.1\,L}。\\
\bkqokindent 1\,L と、 0.1\,L が 3こ。

\kotae 1.3\,L
\end{minipage}\hfill
\begin{minipage}[c]{0.32\linewidth}\centering
\begin{tikzpicture}[x=1mm, y=1mm]
  \bkqomasu{(0,0)}{14}{22}{10}\bkqomasu{(18,0)}{14}{22}{3}
  \node[font=\sffamily\bfseries\footnotesize] at (0,-3.5) {1\,L};
  \node[font=\sffamily\bfseries\footnotesize] at (18,-3.5) {1\,L};
\end{tikzpicture}
\end{minipage}
\end{reidai}

\begin{jibun}
□に あてはまる \R{数}{かず}を かこう。\par\vspace{2pt}
\noindent\hspace*{1\zw}（1）2\,L 5\,dL ＝ \bkkakikomi\,L\hspace{3\zw}（2）7\,dL ＝ \bkkakikomi\,L
\bkjans{（1）2.5\quad（2）0.7}
\end{jibun}

\begin{reidai}[かずの \bkqofrac{1}{2}・\bkqofrac{1}{4}]
いちごが 12こ あります。 （1）12この \bkqofrac{1}{2} は なんこ？\hspace{1\zw}（2）12この \bkqofrac{1}{4} は なんこ？

\kangaekata （1）2つに \R{同}{おな}じ かずずつ \R{分}{わ}ける。\\
\bkqokindent\hspace*{2\zw}\mbox{12 ÷ 2 ＝ 6}\\
\bkqokindent（2）4つに \R{同}{おな}じ かずずつ \R{分}{わ}ける。\\
\bkqokindent\hspace*{2\zw}\mbox{12 ÷ 4 ＝ 3}

\kotae （1）6こ\hspace{2\zw}（2）3こ
\end{reidai}

\begin{jibun}
（1）10この \bkqofrac{1}{2} は なんこ？\hfill\bkkakikomi\enspace こ\par\vspace{2pt}
（2）16この \bkqofrac{1}{4} は なんこ？\hfill\bkkakikomi\enspace こ
\bkjans{（1）\mbox{10 ÷ 2 ＝ 5}\quad 5こ\quad（2）\mbox{16 ÷ 4 ＝ 4}\quad 4こ}
\end{jibun}

\bkfukidashi{delta}{\bkqofrac{1}{2} は 「2つに \R{同}{おな}じに \R{分}{わ}けた 1つ」。 だから かずの \bkqofrac{1}{2} は 「÷ 2」、 \bkqofrac{1}{4} は 「÷ 4」で もとめられるんだ。}

\begin{reidai}[\R{小数}{しょう|すう}と \R{分数}{ぶん|すう}]
\noindent\begin{minipage}[c]{0.56\linewidth}
1\,L の \bkqofrac{1}{2} は、 \R{何}{なん}\,L？ \R{小数}{しょう|すう}で こたえよう。

\kangaekata 1\,L ますの \R{目}{め}もりは 10こ。 その \R{半分}{はん|ぶん}は 5こ。 0.1\,L が 5こで 0.5\,L。

\kotae 0.5\,L
\end{minipage}\hfill
\begin{minipage}[c]{0.42\linewidth}\centering
\begin{tikzpicture}[x=1mm, y=1mm]
  \bkqomasu{(0,0)}{16}{26}{5}
  \draw[bkTheme, line width=0.8pt] (9,13) -- (13,13);
  \node[anchor=west, font=\sffamily\bfseries\footnotesize, text=bkTheme, align=left] at (13.5,13) {\R{半分}{はん|ぶん}\\（\R{目}{め}もり 5こ）};
  \node[font=\sffamily\bfseries\footnotesize] at (0,-3.5) {1\,L ます};
\end{tikzpicture}
\end{minipage}
\end{reidai}

\begin{jibun}
（1）10\,cm の テープの \bkqofrac{1}{2} は、 なん\,cm？\hfill\bkkakikomi\,cm\par\vspace{2pt}
（2）0.5 を \R{分数}{ぶん|すう}で かこう。\hfill\bkkakikomi
\bkjans{（1）\mbox{10 ÷ 2 ＝ 5}\quad 5\,cm\quad（2）\bkqofrac{1}{2}}
\end{jibun}

%% ---------------- p.166 まちがえやすいところ ----------------
\bksection{tsumazuki}{まちがえやすい ところ}

\begin{tsumazuki}{\R{大}{おお}きさが ちがう}
\noindent\begin{minipage}[c]{0.76\linewidth}
これは \bkqofrac{1}{2}？\hspace{1\zw}\bkbatsu\enspace はい\hspace{1\zw}\bkmaru\enspace いいえ\\
\R{同}{おな}じ \R{大}{おお}きさで ないから、 \bkqofrac{1}{2} では ないよ。
\end{minipage}\hfill
\begin{minipage}[c]{0.22\linewidth}\centering
\begin{tikzpicture}[x=1mm, y=1mm]
  \fill[bkYellowBg] (0,0) circle (10);
  \begin{scope}\clip (0,0) circle (10); \fill[bkRed!45] (4,-11) rectangle (11,11);\end{scope}
  \draw[bkInk!80, line width=0.8pt] (4,-9.17) -- (4,9.17);
  \draw[bkYellow!60!black, line width=1.3pt] (0,0) circle (10);
\end{tikzpicture}
\end{minipage}
\end{tsumazuki}

\begin{tsumazuki}{4は 2より \R{大}{おお}きいから…}
\noindent\begin{minipage}[c]{0.7\linewidth}
\R{大}{おお}きいのは？\hspace{1\zw}\bkbatsu\enspace\bkqofrac{1}{4}\hspace{1.5\zw}\bkmaru\enspace\bkqofrac{1}{2}\\
たくさんに \R{分}{わ}けるほど、 1つ\R{分}{ぶん}は \R{小}{ちい}さく なるよ。
\end{minipage}\hfill
\begin{minipage}[c]{0.28\linewidth}\centering
\begin{tikzpicture}[x=1mm, y=1mm]
  \bkpizza{(0,0)}{8}{2}{1}{bkRed!45}
  \bkpizza{(20,0)}{8}{4}{1}{bkRed!45}
  \node[font=\sffamily\bfseries\footnotesize] at (0,-13) {\bkqofrac{1}{2}};
  \node[font=\sffamily\bfseries\footnotesize] at (20,-13) {\bkqofrac{1}{4}};
\end{tikzpicture}
\end{minipage}
\end{tsumazuki}

\begin{tsumazuki}{\R{小数点}{しょう|すう|てん}を わすれる}
3\,cm 4\,mm は？\hspace{1\zw}\bkbatsu\enspace 34\,cm\hspace{2\zw}\bkmaru\enspace 3.4\,cm\par
「.」を わすれると、 ずっと \R{長}{なが}く なって しまうよ。
\end{tsumazuki}

\begin{tsumazuki}{0.9の つぎ}
0.1ずつ ふやすと、 0.7、 0.8、 0.9、 その つぎは？\\
\bkbatsu\enspace 0.10\hspace{3\zw}\bkmaru\enspace 1\par
0.1が 10こ あつまると 1に なるよ。 1-2の 「10こ あつまると、 1つ \R{左}{ひだり}の くらいへ すすむ」と \R{同}{おな}じだね。
\end{tsumazuki}

\begin{column}{\R{分数}{ぶん|すう}は とても \R{古}{ふる}い}
\mbox{3500\R{年}{ねん}}より むかしの エジプトの \R{紙}{かみ}（パピルス）に、 パンを \mbox{10\R{人}{にん}}で \R{同}{おな}じに \R{分}{わ}ける \mbox{\R{計算}{けい|さん}}が かかれて います。 \mbox{\R{分数}{ぶん|すう}}は、 「\R{同}{おな}じに \R{分}{わ}けたい」から \R{生}{う}まれた \R{数}{かず}なのです。
\end{column}

\clearpage
%% ---------------- p.167 つかわれているところ ----------------
\begin{butsuri}
1より \R{小}{ちい}さい はんぱも、 \mbox{\R{小数}{しょう|すう}}や \mbox{\R{分数}{ぶん|すう}}で はかって います。
\begin{bkfigure}
\begin{tikzpicture}[x=1mm, y=1mm]
  % 体温計
  \fill[white, rounded corners=3mm] (0,0) rectangle (38,34);
  \fill[white, draw=bkGray, line width=0.8pt, rounded corners=3mm] (3,14) rectangle (35,24);
  \fill[bkGray!50] (3.6,17) rectangle (7,21);
  \fill[bkGreen!25, rounded corners=0.8mm] (12,15.8) rectangle (33,22.2);
  \node[anchor=east, font=\sffamily\bfseries\small, text=bkInk] at (32.5,19) {36.5\,℃};
  \node[font=\sffamily\bfseries\footnotesize] at (19,6) {\R{体温計}{たい|おん|けい}};
  % ガソリンの メーター
  \begin{scope}[shift={(42,0)}]
  \fill[white, rounded corners=3mm] (0,0) rectangle (38,34);
  \draw[bkInk!80, line width=1.2pt] (7,10.5) arc (180:0:12);
  \foreach \a in {180,135,90,45,0} {\draw[bkInk!80, line width=0.9pt] ($(19,10.5)+(\a:12)$) -- ($(19,10.5)+(\a:9.6)$);}
  \node[font=\sffamily\bfseries\footnotesize] at (7,7.3) {E};
  \node[font=\sffamily\bfseries\footnotesize] at (31,7.3) {F};
  \node[font=\sffamily\bfseries\footnotesize] at (19,28.6) {\bkqofrac{1}{2}};
  \draw[bkRed, line width=1.4pt, line cap=round] (19,10.5) -- ++(90:10);
  \fill[bkInk] (19,10.5) circle (1.4);
  \node[font=\sffamily\bfseries\footnotesize] at (19,4.5) {ガソリン};
  \end{scope}
  % かけっこの タイム
  \begin{scope}[shift={(84,0)}]
  \fill[white, rounded corners=3mm] (0,0) rectangle (38,34);
  \fill[bkGray!60] (17.5,27) rectangle (20.5,30);
  \filldraw[fill=bkBlue!15, draw=bkBlue!60!black, line width=1pt] (19,17.5) circle (9.5);
  \fill[white, rounded corners=0.8mm] (11.5,14.3) rectangle (26.5,20.7);
  \node[font=\sffamily\bfseries\small, text=bkInk] at (19,17.5) {9.8};
  \node[font=\sffamily\bfseries\footnotesize] at (19,4.5) {タイム（\R{秒}{びょう}）};
  \end{scope}
\end{tikzpicture}
\end{bkfigure}
\begin{itemize}[label={\color{bkPurple}●}, leftmargin=1.4\zw, itemsep=0pt, topsep=0pt]
\item \mbox{\R{体温計}{たい|おん|けい}}（36.5\,℃）や かけっこの タイム（\mbox{9.8\R{秒}{びょう}}）は、 0.1ずつ はかります。
\item ガソリンの メーターの はりが まん\R{中}{なか}なら、 のこりは \bkqofrac{1}{2}（\mbox{\R{半分}{はん|ぶん}}）。
\item \mbox{\R{工場}{こう|じょう}}では、 \mbox{\R{部品}{ぶ|ひん}}の \R{大}{おお}きさを 0.1\,mm より こまかく はかって、 そろえて \R{作}{つく}ります。 0.1\,mm は、 かみの \R{毛}{け} \mbox{1\R{本}{ぽん}}の \R{太}{ふと}さ ほどです。
\end{itemize}
\end{butsuri}

\bkfukidashi{delta}{はんぱまで はかれると、 ぴったりの ものが \R{作}{つく}れるね。 「はかる」の \mbox{\R{路線}{ろ|せん}}は、 \mbox{\R{大学}{だい|がく}}の \mbox{\R{計測}{けい|そく}}へ つづくよ！}

\begin{bkfigure}
\begin{tikzpicture}[x=1mm, y=1mm]
  \draw[bkRoute7, line width=3pt, line cap=round] (10,0) -- (128,0);
  \node[anchor=west, inner sep=0pt] at (0,12) {\bkrouteicon[1.4em]{7}\hspace{0.3em}{\sffamily\bfseries\small\color{bkRoute7!80!black}「はかる」の \mbox{\R{路線}{ろ|せん}}}};
  \foreach \x/\c/\t/\s in {10/bkVol1/{1-9}/{かさ・おもさ},40/bkVol1/{1-11}/{\R{目}{め}もり},70/bkVol1/{1-16}/{\R{小数}{しょう|すう}},100/bkVol2/{2-6}/{ずれ},128/bkVol3/{3-14}/{ばらつき}} {
    \fill[white, draw=\c, line width=1.6pt] (\x,0) circle (2.6);
    \node[anchor=base, font=\sffamily\bfseries\footnotesize, text=\c!85!black] at (\x,-6.8) {\t};
    \node[anchor=base, font=\sffamily\footnotesize, text=bkInk] at (\x,-13.6) {\s};}
  \node[fill=bkTheme, text=white, rounded corners=1mm, font=\sffamily\bfseries\fontsize{8}{9}\selectfont, inner sep=1.5pt] at (70,5.5) {いま ここ};
\end{tikzpicture}
\end{bkfigure}

\clearpage
%% ---------------- p.168〜169 もんだい ----------------
\bkendproblems

\bkproblemgroup{きほん}{1}

\bkproblem □に あてはまる \R{数}{かず}を かきましょう。
\par\vspace{2pt}
\noindent\hspace*{1\zw}\begin{tblr}{colspec={Q[l,m,wd=72mm]Q[l,m,wd=60mm]}, colsep=0pt, rowsep=3pt}
（1）2\,cm 7\,mm ＝ \bkkakikomi\,cm & （2）5\,mm ＝ \bkkakikomi\,cm \\
（3）1.5は 0.1が \bkkakikomi こ & （4）0.1が 10こで \bkkakikomi \\
（5）3\,L 2\,dL ＝ \bkkakikomi\,L & （6）4\,dL ＝ \bkkakikomi\,L \\
\end{tblr}
\bkpans{（1）2.7\quad（2）0.5\quad（3）15\quad（4）1\quad（5）3.2\quad（6）0.4}

\bkproblem \R{水}{みず}の かさは \R{何}{なん}\,dL？\hfill
\begin{tikzpicture}[x=1mm, y=1mm, baseline=(current bounding box.center)]
  \bkqomasu{(0,0)}{12}{18}{10}\bkqomasu{(15,0)}{12}{18}{6}
\end{tikzpicture}\hspace{6mm}{\sffamily\bfseries\small\color{bkOrange}\bkkotaelabel}\enspace\bkkakikomi[24mm]
\bkpans{1.6\,dL}

\bkproblem ア、 イ、 ウの \R{目}{め}もりは いくつ？
\begin{bkfigure}
\begin{tikzpicture}[x=1mm, y=1mm]
  \bkqonum{(0,0)}{55}{2}
  \bkqomark{(11,0)}{ア}{bkTheme}
  \bkqomark{(49.5,0)}{イ}{bkTheme}
  \bkqomark{(77,0)}{ウ}{bkTheme}
\end{tikzpicture}
\end{bkfigure}
\noindent\hspace*{\fill}ア\enspace\bkkakikomi[14mm]\hspace{1.5\zw}イ\enspace\bkkakikomi[14mm]\hspace{1.5\zw}ウ\enspace\bkkakikomi[14mm]
\bkpans{\mbox{ア 0.2}、 \mbox{イ 0.9}、 \mbox{ウ 1.4}}

\bkproblem \R{色}{いろ}の ところは、 もとの \R{大}{おお}きさの どれだけ？
\begin{bkfigure}
\begin{tikzpicture}[x=1mm, y=1mm]
  \node[font=\sffamily\bfseries\small] at (0,0) {（1）};
  \bkqotape{(8,-3.5)}{36}{2}{1}{bkTheme!60}
  \node[anchor=west, font=\sffamily\bfseries\small] at (46,0) {\bkkakikomi[14mm]};
  \node[font=\sffamily\bfseries\small] at (70,0) {（2）};
  \fill[bkTheme!60] (78,0) rectangle (86,8);
  \draw[bkInk!80, line width=0.8pt] (78,-8) rectangle (94,8) (86,-8) -- (86,8) (78,0) -- (94,0);
  \node[anchor=west, font=\sffamily\bfseries\small] at (96,0) {\bkkakikomi[14mm]};
\end{tikzpicture}
\end{bkfigure}
\bkpans{（1）\bkqofrac{1}{2}\quad（2）\bkqofrac{1}{4}}

\bkproblem \R{大}{おお}きい ほうに ○を つけよう。\par\vspace{2pt}
\noindent\hspace*{1\zw}（1）0.6\enspace と\enspace 1.1\hspace{3\zw}（2）2.4\enspace と\enspace 2.8
\bkpans{（1）1.1\quad（2）2.8}

\bkproblemgroup{チャレンジ}{1}

\bkproblem えんぴつの \R{長}{なが}さは \R{何}{なん}\,cm？\hfill{\sffamily\bfseries\small\color{bkOrange}\bkkotaelabel}\enspace\bkkakikomi[24mm]
\begin{bkfigure}
\begin{tikzpicture}[x=1mm, y=1mm]
  \bkpencil{(2.2,4)}{70.8}{bkRoute2}
  \bkqoruler{(0,0)}{8}
\end{tikzpicture}
\end{bkfigure}
\bkpans{7.3\,cm}

\bkproblem \R{紙}{かみ}テープを \mbox{\R{半分}{はん|ぶん}}に \R{切}{き}って、 その \mbox{1\R{本}{ぽん}}を また \mbox{\R{半分}{はん|ぶん}}に \R{切}{き}りました。 いちばん みじかい テープは、 もとの テープの どれだけ？\hfill{\sffamily\bfseries\small\color{bkOrange}\bkkotaelabel}\enspace\bkkakikomi[20mm]
\bkpans{\bkqofrac{1}{4}}

\bkproblem クッキーが 20まい あります。 （1）20まいの \bkqofrac{1}{2} は なんまい？\hspace{1\zw}（2）20まいの \bkqofrac{1}{4} は なんまい？
\bkshikiran
\bkshikiran
\bkpans{（1）\mbox{20 ÷ 2 ＝ 10}\quad 10まい\quad（2）\mbox{20 ÷ 4 ＝ 5}\quad 5まい}

\begin{furikaeri}
\begin{checklist}
\item はんぱを \mbox{\R{小数}{しょう|すう}}で あらわせる。\hfill{\footnotesize\sffamily\color{bkTheme}→\pageref{bk:1-16:shousuu}ページ}
\item \mbox{\R{数直線}{すう|ちょく|せん}}で \mbox{\R{小数}{しょう|すう}}が よめる。\hfill{\footnotesize\sffamily\color{bkTheme}→\pageref{bk:1-16:suuchokusen}ページ}
\item \mbox{\R{半分}{はん|ぶん}}を \mbox{\R{分数}{ぶん|すう}}で あらわせる。\hfill{\footnotesize\sffamily\color{bkTheme}→\pageref{bk:1-16:bunsuu}ページ}
\end{checklist}
チェックが つかなかったら、 →の ページに もどろう。
\end{furikaeri}

\bkfukidashi{sora}{おうちの \mbox{\R{体温計}{たい|おん|けい}}や ペットボトルで、 \mbox{\R{小数}{しょう|すう}}や \mbox{\R{分数}{ぶん|すう}}を さがして みよう！}
\clearpage
